\documentclass[12pt]{article}

\usepackage[margin=1.25in]{geometry}
\usepackage{setspace}
\usepackage{hyperref}
\usepackage{footnote}
\usepackage{booktabs}
\usepackage[T1]{fontenc}

\hypersetup{
    colorlinks=true,
    linkcolor=blue,
    citecolor=blue,
    urlcolor=blue
}

%\doublespacing

\title{Two Threads Pulled:\\[0.5em]
Madison's Convention Notes on the ``Emit Bills'' Debate\\
and Justice Field's Dissents in the Legal Tender Cases}
\author{Notes on Bancroft, Primary Sources, and the Logic of Constitutional Commitment}
\date{\today}

\begin{document}

\maketitle
\thispagestyle{empty}

\newpage
\tableofcontents
\newpage

%-----------------------------------------------------------------------
\section{Introduction: Two Interpretive Puzzles}
%-----------------------------------------------------------------------

The previous document on Bancroft's \textit{Plea for the Constitution}
closed with two threads worth pulling. The first concerns the precision
of Bancroft's historical claim that the Constitutional Convention of 1787
``shut and barred the doors against paper money'' by deliberate
supermajority vote---a claim that can now be checked against the
primary source Bancroft himself cites, James Madison's convention
minutes for Thursday, August~16, 1787. The second concerns Justice
Stephen Field's dissents in the Legal Tender Cases, which Bancroft
treats as vindicated by history but which are worth examining closely
for their own constitutional logic, particularly the argument that the
principle ``where a power is not expressly forbidden it may be
exercised'' would fundamentally transform the character of the American
constitutional order.

Both threads converge on a single hermeneutical crux: the text of Madison's
own footnote to his record of the August~16 vote, which the \textit{Juilliard}
majority turned against Bancroft's reading to argue that striking out
``emit bills'' had deliberately left the legal tender question unresolved.
This makes the episode more genuinely ambiguous than either Bancroft or
his opponents admitted---and connects it, by a deep structural logic, to the
ambiguity we have traced in the 5-20 bond legislation.

%-----------------------------------------------------------------------
\section{The Primary Source: Madison's Notes for August 16, 1787}
%-----------------------------------------------------------------------

The relevant passage in Madison's notes is brief---under twenty lines---
but its brevity conceals a remarkable range of positions and a carefully
recorded footnote that became the pivot of later constitutional
controversy. The full account of the ``emit bills'' debate is worth
reconstructing in detail.

\subsection{The Motion and Its Framing}

The committee of detail's draft of what would become Article~I,
Section~8 read, in the relevant part: the Congress shall have power
``to borrow money and emit bills on the credit of the United States.''
Gouverneur Morris moved to strike out the words ``and emit bills on
the credit of the United States.'' His reasoning was lapidary:
``If the United States had credit such bills would be unnecessary:
if they had not, unjust and useless.''

This was not a merely prudential argument about the dangers of
depreciation; it was a logical syllogism. Good credit made paper
notes redundant---creditors would lend on ordinary terms without
the coercive device of legal tender. Bad credit made them a fraud---
forcing creditors to accept a depreciating instrument at face value
was an act of dispossession in statutory disguise. Butler seconded the
motion.

\subsection{Madison's Initial Position: A Middle Path}

Madison's first response was not to support Morris's motion but to
propose a softer alternative: ``will it not be sufficient to prohibit
the making them a tender? This will remove the temptation to emit them
with unjust views. And promissory notes in that shape may in some
emergencies be best.''

This is a crucial moment. Madison in August 1787 initially favored
retaining the power to emit bills while stripping the legal tender
quality---precisely the structure that the Legal Tender Act of 1862
would later adopt for interest payments on the 5-20s (coin mandated)
while leaving the principal currency question open. The instrument he
envisaged was a government bill that circulated by voluntary market
acceptance rather than legal compulsion---close to what modern central
banking terminology would call ``high-powered money'' without legal
tender status.

Morris's reply sharpened the analysis: striking out the words entirely
would leave room still for ``notes of a responsible minister''---
a phrase that bears emphasis---``which will do all the good without
the mischief.'' The ``monied interest'' would oppose any constitution
that retained the paper emission power. Here is the reputational logic
again: if the power existed in statutory form, investors would price
it into bond yields; if it was credibly committed away, borrowing
costs would fall. This is precisely the Rogoff (1985) logic of
institutional commitment as a substitute for the discretion of the
current officeholder.

\subsection{The Range of Positions}

The subsequent debate divided the delegates into several camps, recorded
with unusual clarity in Madison's minutes:

\textit{For striking out entirely without any prohibition:} Gorham
(Massachusetts) argued that if the words stood, they would ``suggest
and lead to the measure,'' and that ``the power, so far as it will
be necessary or safe, is involved in that of borrowing.'' Ellsworth
(Connecticut) called it ``a favorable moment to shut and bar the door
against paper money,'' adding that ``the mischiefs of the various
experiments which had been made, were now fresh in the public mind
and had excited the disgust of all the respectable part of America.''
Wilson (Pennsylvania) thought ``it will have a most salutary influence
on the credit of the U. States to remove the possibility of paper
money.'' Butler (South Carolina) noted that ``paper was a legal tender
in no Country in Europe.'' And Langdon (New Hampshire) delivered the
meeting's most dramatic line: he ``had rather reject the whole plan
than retain the three words.'' Read (Delaware) thought the words,
if not struck, ``would be as alarming as the mark of the Beast in
Revelations.''

\textit{For retaining the power:} Randolph (Virginia) ``notwithstanding
his antipathy to paper money, could not agree to strike out the words,
as he could not foresee all the occasions that might arise.'' Mason
(Virginia) professed ``a mortal hatred to paper money'' but was
``unwilling to tie the hands of the Legislature,'' observing that
``the late war could not have been carried on, had such a prohibition
existed.'' Mercer (Maryland) was more candidly in favor of retaining
the power, arguing that it ``would stamp suspicion on the Government
to deny it a discretion on this point.''

\textit{Voted against the motion but signed the Constitution:}
New Jersey alone voted no, along with Maryland. Crucially, both
Randolph and Mason---who had spoken against the motion---are recorded
as voting with Virginia's affirmative, because of what Madison's
footnote tells us.

\subsection{The Vote and Madison's Footnote}

The motion carried 9--2 (New Hampshire, Massachusetts, Connecticut,
Pennsylvania, Delaware, Virginia, North Carolina, South Carolina,
and Georgia voting aye; New Jersey and Maryland no).

Virginia's affirmative vote was unexpected, given that both Randolph
and Mason had spoken against the motion. Madison's footnote to his
own record of the vote explains: ``This vote in the affirmative by
Virga. was occasioned by the acquiescence of Mr. Madison who became
satisfied that striking out the words would not disable the Govt.
from the use of public notes as far as they could be safe and proper;
and would only cut off the pretext for a paper currency, and
particularly for making the bills a tender either for public or
private debts.''

This footnote is the pivot of the entire subsequent controversy, and it
rewards extremely careful reading.

%-----------------------------------------------------------------------
\section{Bancroft's Account vs.\ Madison's Notes: Where They Agree
and Diverge}
%-----------------------------------------------------------------------

\subsection{Where Bancroft Is Accurate}

Bancroft's account in Part~III of the \textit{Plea} is substantially
accurate in its core claim: the Convention deliberately struck the
``emit bills'' language, did so by a large and determined majority
(9--2), and most delegates understood themselves to be cutting off the
pretext for paper legal tender. The direct quotations Bancroft provides
from the debates---Morris's syllogism, Wilson's ``most salutary
influence,'' Gorham's argument, Mason's ``mortal hatred''---all
appear in Madison's minutes. Bancroft also accurately records
Madison's own footnote explaining the Virginia vote.

On the crucial historical question of intent, Bancroft has the better
of the argument against the Juilliard majority. The delegates who
spoke at greatest length in favor of striking out---Morris, Wilson,
Ellsworth, Langdon---were not merely removing an inconvenient grant;
they were making an affirmative constitutional statement about the
kind of monetary system the new government would have. Gorham's
argument that the power ``so far as it will be necessary or safe, is
involved in that of borrowing'' is the germ of what became the
\textit{McCulloch/Juilliard} majority position---but Gorham was the
most cautious voice among those who wanted to strike, and even he
said only ``so far as it will be necessary or safe,'' not ``without
limit.''

\subsection{The Hermeneutical Ambiguity That Bancroft Downplays}

Yet Madison's footnote is genuinely ambiguous in a way that Bancroft
does not fully confront, and that Justice Gray's majority in
\textit{Juilliard} exploited effectively.

Madison writes that striking out the words would ``not disable the
Govt. from the use of public notes as far as they could be safe and
proper.'' This sentence has two possible readings.

Under Bancroft's reading, ``as far as they could be safe and proper''
means: the government may issue bills that circulate voluntarily, but
may not give them the force of legal tender, which is the unsafe and
improper application. The Convention struck the power to make bills
legal tender, while leaving the power to issue voluntary instruments.

Under Gray's reading in \textit{Juilliard}, the footnote establishes
that Madison \textit{deliberately chose not to prohibit the power
entirely}, because he thought Congress would retain some bill-issuing
capacity from its borrowing power. The question of legal tender status
was therefore not definitively resolved by the ``emit bills'' deletion.
Gray writes: ``But he has not explained why he thought that striking
out the words `and emit bills' would leave the power to emit bills,
and deny the power to make them a tender in payment of debts.''

This is a genuine textual difficulty. Madison's footnote says the
deletion would cut off the pretext for ``making the bills a tender''---
which implies he thought the tender power was eliminated. But his
initial motion in the debate (``will it not be sufficient to prohibit
the making them a tender?'') had treated the tender quality as
separable from the emission power---and he abandoned that motion in
favor of Morris's more complete deletion. If Madison really believed
the deletion eliminated legal tender but not emission, as his footnote
asserts, then Gray's question---why did he think so?---has no clean
answer from the text alone.

The honest assessment is that Madison himself held the ambiguity in
his own mind. He believed the Convention had cut off the pretext for
legal tender, but the textual mechanism by which that occurred was
implicit (emission power gone, hence tender power gone too) rather
than explicit (a textual prohibition on the federal government's making
paper a legal tender, analogous to the explicit prohibition on the
states in Article~I, Section~10). The absence of that explicit
prohibition is precisely what both the \textit{Knox v.\ Lee} majority
and the \textit{Juilliard} majority hung their arguments on.

\subsection{Bancroft vs.\ Madison on the Convention's Affirmative Design}

Where Bancroft is on firmest ground is in documenting the \textit{spirit}
of the August~16 vote, independent of the precise textual mechanism.
The delegates who spoke most forcefully---and who voted aye---were not
engaged in a technical act of drafting economy. Langdon would have
rejected the entire Constitution to excise three words. Ellsworth
called it ``a favorable moment to shut and bar the door.'' Wilson
described the effect as ``removing the possibility of paper money.''
Read compared the words to the mark of the Beast.

This is not the language of men who thought they were leaving a
discretionary power in the borrowing clause. And Madison's own
footnote---``would only cut off the pretext for a paper currency,
and particularly for making the bills a tender''---uses the word
``only'' as a limiting particle: the sole effect of the deletion was
to eliminate paper legal tender, not to disable government borrowing.
The burden of Bancroft's argument is that the Convention's
demonstrable intent---as expressed both in the recorded debate and
in Madison's own explanatory note---cannot be overridden by the
interpretive technique of finding silence where the Framers evidently
expected prohibition to be implied.

%-----------------------------------------------------------------------
\section{Justice Field's Dissents: A Close Reading}
%-----------------------------------------------------------------------

Field dissented in all three Legal Tender Cases---\textit{Hepburn v.\
Griswold} (1870, though there he was in the majority under Chase),
\textit{Knox v.\ Lee} (1871), and \textit{Juilliard v.\ Greenman}
(1884). His \textit{Knox} dissent is the fullest statement of his
position; in \textit{Juilliard} he simply adhered to the views expressed
in the earlier cases, and the majority noted this without attempting
to answer his arguments on the merits. Field's dissent richly repays
close reading as a piece of constitutional reasoning.

\subsection{The Positive Claim: ``Coin Money'' Is Incompatible with
Paper Legal Tender}

Field opens his substantive constitutional argument with a point that
is more powerful than it initially appears. The power ``to coin money,''
he argues, is not merely in tension with a power to make paper legal
tender; it is structurally incompatible with one. To coin money is
to mould metallic substances having intrinsic value into certain forms
convenient for commerce, and to impress them with the stamp of the
government indicating their value. The point of coinage is precisely
that the stamp attests to an intrinsic value: the coin is worth what
the government says it is worth because the metal in it is worth that.

The legal tender power, by contrast, divorces denomination from
intrinsic value. It says: this piece of paper, which has no intrinsic
value, shall nevertheless discharge a debt of one hundred dollars in
full. This is not a regulation of monetary value; it is a legal
compulsion substituted for monetary value. ``Such a medium,'' Chief
Justice Marshall had written of paper money, ``has been always liable
to considerable fluctuation. Its value is continually changing; and
these changes, often great and sudden, expose individuals to immense
loss, are the sources of ruinous speculations, and destroy all
confidence between man and man.''

Field's argument is that the Framers, in giving Congress the power
to coin money and regulate its value, and in giving the states the
power to make nothing but gold and silver coin a tender, had established
a coherent metallic monetary system in which the legal tender quality
and the intrinsic value of the tender were inseparable. To read the
federal borrowing power as implicitly containing a legal tender power
was to read into the Constitution an instrument that its express
monetary provisions were designed to exclude.

The inhibition on the states to make anything but gold and silver coin
a tender, Field writes, ``must be read in connection with the grant of
the coinage power to Congress. The two provisions taken together
indicate beyond question that the coins which the National government
was to fabricate, and the foreign coins, the valuation of which it
was to regulate, were to consist principally, if not entirely, of
gold and silver.'' The structure of the monetary provisions---the
positive coinage power, the state prohibition, the foreign coin
valuation---forms a system in which paper legal tender has no place.

\subsection{The Negative Claim: ``The Wants of the Government Can
Never Be the Measure of Its Powers''}

Field's second main argument targets the reasoning by which both
\textit{Knox} and \textit{Juilliard} derived the legal tender power
from the war power and the borrowing power. The core of that reasoning,
as Field understood it, was: the Civil War created an urgent need for
funds; paper money was the means Congress chose; therefore, having the
power to conduct war and to borrow money, Congress had the implied power
to make its paper legal tender.

Field objects that this reasoning makes the magnitude of government
want the measure of government power. ``The existence of the war only
increased the urgency of the government for funds. It did not add to
its powers to raise such funds, or change, in any respect, the nature
of those powers or the transactions which they authorized. \ldots\ The
wants of the government can never be the measure of its powers.''

This is one of the most important sentences in the history of American
constitutional law, and it deserves to be stated in its full force.
Under the implied powers doctrine as Bancroft's opponents used it,
extreme necessity generates extreme power: because the Union could not
survive without paper money in 1862, Congress must have had the
constitutional power to issue it. Field's response is that this is
precisely backwards. The enumerated powers doctrine is not an all-weather
principle that applies only in ordinary times and yields in emergencies;
it is a constraint that applies \textit{precisely} in emergencies,
when the temptation to override it is greatest. An emergency that
expands implied powers is an emergency that effectively rewrites the
Constitution without amendment---converting a document of limited
enumerated powers into an instrument of unlimited sovereign discretion.

The force of this argument can be appreciated by comparing it to the
structure of the 5-20 bond episode. Chase argued in 1862 that the
war emergency necessitated legal tender; he later argued in
\textit{Hepburn} that the same emergency had not created constitutional
authority for it. Field's response to both positions is consistent:
emergency necessity is not a constitutional source of power, whether
we are evaluating the original act or evaluating it retrospectively.
The wants of the government in 1862 were real and urgent. They were
also constitutionally irrelevant to the scope of congressional power.

\subsection{The Enumerated-Powers Syllogism: The Bankruptcy Argument}

Field's most precise constitutional argument is a negative inference from
the express grant of bankruptcy power. Article~I, Section~8 gives
Congress the power to establish uniform laws of bankruptcy. This is the
one express provision by which Congress may alter the obligation of
contracts---specifically, by releasing insolvent debtors from their debts
upon surrender of their property.

Field's inference is that this \textit{express} grant is a strong
argument against any \textit{implied} general power in Congress to
interfere with the obligation of contracts. The constitutional structure
is: if Congress needed an express grant to authorize bankruptcy (which
releases debtors from their contracts in a supervised and equitable
proceeding), then Congress certainly needed an express grant to authorize
compulsory partial discharge at a depreciated value (which is what
legal tender for private debts amounts to). The express grant proves
that the Framers understood this class of legislation as requiring
explicit authorization, not implied power.

This argument cuts directly against the \textit{McCulloch} framework
as applied in \textit{Knox} and \textit{Juilliard}. Under
\textit{McCulloch}, necessary and proper means are constitutional even
when not expressly enumerated, provided they are conducive to the
execution of an enumerated power. Field's bankruptcy argument says:
the structure of the Constitution itself tells us that this particular
class of means---interference with contractual obligations---requires
express authorization, because the Framers put in an express grant
for the specific, limited version of it they approved. To read implied
power here is to render the express bankruptcy clause redundant.

\subsection{The Master Argument: ``The Doctrine That Where a Power Is
Not Expressly Forbidden It May Be Exercised''}

Field's culminating argument is the most fundamental, and it is the
one that Bancroft endorses most emphatically. The Juilliard majority's
reasoning, stripped to its logical structure, was: the power to make
paper legal tender is not expressly withheld from Congress by the
Constitution; therefore, being one of the powers belonging to sovereignty
in other civilized nations, it is included in the implied powers available
to Congress.

Field's response: ``The doctrine that where a power is not expressly
forbidden it may be exercised, would change the whole character of
our government.''

The logic of this response requires unpacking. The Tenth Amendment
reserves to the states or the people all powers not delegated to the
federal government. This means the federal government's powers are
defined by positive delegation, not by residual exception. A
constitutional design of enumerated powers says: the government may
do X, Y, and Z (and what is necessary and proper to execute them).
It does not say: the government may do anything except A, B, and C.
The latter design characterizes a government of general powers with
specific exceptions, which is the design of state constitutions, not
the federal constitution.

Gray's majority opinion in \textit{Juilliard}, by finding legal tender
power in the proposition that it is ``not expressly withheld,'' had
implicitly converted the federal constitution from a grant-of-powers
document into an exception-from-powers document. This is not a minor
interpretive shift; it is a transformation of the constitutional
structure. On Gray's logic, every power belonging to European sovereignties
that Congress has not been expressly forbidden is available to Congress.
That is not a constitution of limited enumerated powers.

Field had delivered the same argument's logical culmination in his
dissent in \textit{Knox}: ``I do not yield to any one in honoring and
reverencing the noble and patriotic men who were in the councils of the
nation during the terrible struggle with the rebellion. \ldots But I do
not admit that a blind approval of every measure which they may have
thought essential to put down the rebellion is any evidence of loyalty
to the country. The only loyalty which I can admit consists in obedience
to the Constitution and laws made in pursuance of it.''

This parting shot in the Legal Tender dissents is not mere rhetoric. It
addresses the specific argument that had been used against opponents of
the Legal Tender Acts in 1862 and thereafter: that to question the
constitutional power was to be lukewarm about the Union cause. Field's
response is constitutionally principled. Loyalty to the Constitution is
not the same as loyalty to every emergency measure taken in its defense;
if it were, the constitution would dissolve in any sufficiently grave
emergency---precisely the outcome a written constitution of limited
powers is designed to prevent.

%-----------------------------------------------------------------------
\section{The Interpretive Logic Puzzle: What Did the Convention Do?}
%-----------------------------------------------------------------------

The two threads we have pulled converge on what might be called the
interpretive logic puzzle of the ``emit bills'' deletion. Both Bancroft
and Field claim that the Convention clearly prohibited paper legal
tender. Both the \textit{Knox} majority and the \textit{Juilliard}
majority claim that the Convention left the question unresolved. The
puzzle is: which reading of Madison's notes is correct?

\subsection{The Four Positions}

The positions map neatly onto a two-by-two matrix defined by two
questions: (1) Did the Convention intend to prohibit paper emission,
or merely paper legal tender? (2) Did the Convention successfully
achieve its intent through the textual mechanism of deleting ``emit
bills''?

\textit{Bancroft/Field position:} The Convention intended to prohibit
paper legal tender (and paper emission generally); the deletion
successfully achieved this, and any residual emission power in the
borrowing clause did not include the legal tender quality.

\textit{Gorham/McCulloch/Knox majority position:} The Convention
intended to remove an affirmative grant while leaving open what the
borrowing power implied; the deletion was agnostic on legal tender;
therefore, the implied power question is determined by the borrowing
power, read in light of \textit{McCulloch}.

\textit{Madison's own position (as stated in footnote):} The
Convention intended to cut off the pretext for paper legal tender
specifically; the deletion achieved this by removing the affirmative
basis for such a claim; but Madison did not expressly prohibit paper
legal tender, because he believed---perhaps overoptimistically---that
removing the affirmative grant was sufficient.

\textit{Gray/Juilliard majority position:} Madison's footnote shows
only that Madison personally interpreted the deletion as cutting off
the legal tender pretext; it does not show that other delegates shared
this understanding; and the Convention's failure to include an express
prohibition analogous to Article~I, Section~10's prohibition on the
states is the operative textual datum.

\subsection{The Structural Asymmetry}

What is most striking about this debate is a structural asymmetry
that neither side fully acknowledges. The Convention \textit{was}
willing to write explicit prohibition when it meant it: Article~I,
Section~10 explicitly forbids states to make anything but gold and
silver coin a tender. No comparable provision applies to the federal
government. Bancroft argues that the omission of the ``emit bills''
power was an implicit equivalent prohibition. Gray argues that the
absence of an explicit parallel prohibition is the decisive datum.

Both arguments are available from the text; neither is conclusive.
This is, in the final analysis, an instance of the same structural
problem we have traced throughout this series: the absence of explicit
language creates interpretive ambiguity that is resolved by political
outcomes rather than textual logic. The Convention of 1787, like the
Congress of 1862, left a question formally open that its participants
believed they had closed through implication. The subsequent resolution
of both questions---in 1871 and 1884 for the constitutional question,
in 1869 for the 5-20 principal question---depended on who held
political power when the bill came due.

%-----------------------------------------------------------------------
\section{Field's Dissents and the Commitment Theory Connection}
%-----------------------------------------------------------------------

Field's dissents can be read, through the lens of modern commitment
theory, as making a point about constitutional design that the legal
tender debate brings into sharp relief.

The core commitment problem in monetary policy is that a government
always has a short-run incentive to inflate (or in the nineteenth
century context, to debase or depreciate its obligations), even when
long-run credibility requires restraint. The solution, under the
Kydland-Prescott framework, is to install institutional constraints
that make the inflationary option costly or unavailable---to convert
the problem from one of discretion (where the optimum is time-inconsistent)
to one of rules (where commitment is enforced by the constraint itself).

The Constitutional Convention of 1787, on Bancroft's and Field's
reading, was attempting to do exactly this: to write the time-inconsistency
constraint into the constitutional text by eliminating the federal
legal tender power. The commitment was to be self-enforcing because
it was constitutional---not merely statutory (like the Public Credit
Act of 1869, which could be and in practice was later modified) or
reputational (like Chase's implicit gold-repayment promise on the
5-20 principal, which dissolved in the 1867--68 political fight).

What Field understood, and Bancroft articulated in historical terms,
is that the \textit{Knox}/\textit{Juilliard} line of decisions had
retroactively converted the constitutional commitment from a hard
constraint to a soft one. If the legal tender power exists whenever
the government's wants are urgent enough---whenever a sufficiently
grave emergency demands the ``medicine of the constitution''---then
there is no hard constraint at all. The commitment is endogenous:
the government retains legal tender as an implied reserve power to
be deployed precisely when the temptation to deploy it is greatest,
which is precisely when a hard constraint would bind.

Field's sentence---``the wants of the government can never be the
measure of its powers''---is the constitutional formulation of the
Kydland-Prescott critique of discretion. A rule that suspends itself
whenever rule-following is costly is not a rule; it is discretion
with extra steps. A constitutional prohibition that yields when the
government urgently needs it to yield is not a prohibition; it is a
default position subject to override.

\subsection{The Residual Ambiguity: When Hard Constraints Breed
Harder Crises}

Yet the opposing position---the one held by Mason, Randolph, Gorham,
and ultimately by the \textit{Knox}/\textit{Juilliard} majority---also
has a commitment-theoretic dimension that Field does not adequately
address. If the constitutional constraint is genuinely hard, it may
produce catastrophic discontinuities: a government that cannot borrow
by any means in a genuine military emergency may lose the war, which
is a worse outcome than the monetary depreciation that legal tender
produces. Mason's argument at Philadelphia---``the late war could not
have been carried on, had such a prohibition existed''---is exactly
this emergency-flexibility argument.

The response Field would give (and does give implicitly) is that the
flexibility to issue legal tender in genuine emergencies should be
obtained through constitutional amendment, not through judicial
interpretation. If the Civil War revealed that the constitution as
written was too rigid, the remedy was to amend it. Instead, the
Court amended it by judicial construction, transforming a hard
constraint into a soft one without the democratic process the
amendment procedure requires.

This is also, in miniature, the story of the 5-20 bond principal
ambiguity. The 37th Congress in 1862 had the option of writing an
explicit gold-payment clause into the principal repayment provision,
but did not---because doing so would have cost the votes needed for
passage. The flexibility was preserved through statutory silence,
at the cost of six years of political uncertainty resolved only
by the Public Credit Act of 1869. Field's logic would say: if
the Congress needed the flexibility, it should have written the
greenback-principal option explicitly and borne the political cost;
silence-as-ambiguity is the institutionalization of what Madison
called ``the epidemic malady''---the temptation to delay unresolvable
disagreements by passing them forward.

%-----------------------------------------------------------------------
\section{Synthesis: What the Two Threads Reveal}
%-----------------------------------------------------------------------

Pulling the threads together, the comparison of Madison's notes with
Bancroft's account, and the close reading of Field's dissents, produce
several conclusions.

On the convention, Bancroft is substantially right about the intent
and the facts, but the Gray majority correctly identified a structural
lacuna: the absence of an explicit federal paper legal tender prohibition
analogous to Article~I, Section~10. Madison's footnote is the key
primary source, and it is genuinely ambiguous between Bancroft's
reading (the deletion cut off the legal tender pretext) and Gray's
reading (the deletion left open the question of what the borrowing
power implied). What the primary source does establish, against the
\textit{Juilliard} majority, is that Madison himself at the time
believed the deletion achieved the anti-legal-tender purpose---which
is powerful evidence of original intent, even if the textual
mechanism was imperfect.

On Field's dissents, they are more sophisticated as constitutional
arguments than their reputation suggests. The bankruptcy inference,
the incompatibility-of-coinage argument, and the enumerated-powers
master argument are all serious pieces of constitutional reasoning.
The most important sentence---``the doctrine that where a power is
not expressly forbidden it may be exercised, would change the whole
character of our government''---correctly identifies the constitutional
transformation that \textit{Juilliard} accomplished: a shift from a
constitution of enumerated powers to a constitution of residual
sovereign powers, defined by exception rather than delegation.

On the connection to the 5-20 bond ambiguity: the two ambiguities are
instances of the same underlying political economy problem---the
difficulty of making credible commitments about money through
institutional mechanisms that are themselves endogenous to the
political process. The Convention created an imperfect constitutional
hard constraint; the 37th Congress created a statutory soft commitment;
and the Legal Tender Cases converted the constitutional constraint into
something closer to the statutory arrangement. In each case, the
resolution depended on political outcomes (Grant's election, Grant's
court appointments) rather than on the logic of the original text.

Madison in 1787 and Field in 1871 and 1884 both understood the same
structural point, though in different registers: that a commitment is
only as durable as the institutional mechanism that enforces it, and
that a commitment enforced only by implication---only by the good faith
of future interpreters---is not much of a commitment at all.

%-----------------------------------------------------------------------
\section*{Appendix: Key Passages in Madison's August 16, 1787 Notes}
%-----------------------------------------------------------------------

The following passages are reproduced from Madison's convention minutes
(Max Farrand, ed., \textit{Records of the Federal Convention of 1787},
Vol.~2), as transcribed at \url{consource.org}:

\begin{quotation}
\noindent Mr.\ Govr.\ Morris moved to strike out ``and emit bills on
the credit of the U. States''---If the United States had credit such
bills would be unnecessary: if they had not unjust and useless.

Mr.\ Butler 2ds.\ the motion.

Mr.\ Madison, will it not be sufficient to prohibit the making them
a tender? This will remove the temptation to emit them with unjust
views. And promissory notes in that shape may in some emergencies
be best.

Mr.\ Govr.\ Morris.\ striking out the words will leave room still for
notes of a responsible minister which will do all the good without
the mischief. The Monied interest will oppose the plan of Government,
if paper emissions be not prohibited.

Mr.\ Ghorum was for striking out, without inserting any prohibition.
if the words stand they may suggest and lead to the measure. \ldots\
The power as far as it will be necessary or safe, is involved in that
of borrowing.

Col.\ Mason had doubts on the subject. Congs. he thought would not
have the power unless it were expressed. Though he had a mortal
hatred to paper money, yet as he could not foresee all emergences,
he was unwilling to tie the hands of the Legislature.

Mr.\ Ellsworth thought this a favorable moment to shut and bar the
door against paper money.

Mr.\ Wilson. It will have a most salutary influence on the credit of
the U. States to remove the possibility of paper money.

Mr.\ Butler. paper was a legal tender in no Country in Europe.

Mr.\ Read, thought the words, if not struck out, would be as alarming
as the mark of the Beast in Revelations.

Mr.\ Langdon had rather reject the whole plan than retain the three
words.

On the motion for striking out N. H. ay.\ Mas.\ ay.\ Ct.\ ay.\ N-J.
no.\ Pa.\ ay.\ Del.\ ay.\ Md.\ no.\ Va.\ ay. N. C. ay.\ S. C. ay.\
Geo.\ ay.\ [Ayes~--~9; noes~--~2.]
\end{quotation}

Madison's footnote to Virginia's affirmative vote:

\begin{quotation}
\noindent This vote in the affirmative by Virga. was occasioned by the
acquiescence of Mr.\ Madison who became satisfied that striking out the
words would not disable the Govt.\ from the use of public notes as far
as they could be safe \& proper; \& would only cut off the pretext for
a paper currency and particularly for making the bills a tender either
for public or private debts.
\end{quotation}

%-----------------------------------------------------------------------
\section*{References}
%-----------------------------------------------------------------------

\begin{description}
\setlength\itemsep{0.5em}

\item Bancroft, George. 1886. \textit{A Plea for the Constitution of
the United States of America Wounded in the House of Its Guardians}.
New York: Harper \& Brothers. Full text at
\texttt{constitution.org/2-Authors/gb/gb-plea.htm}.

\item Chase, Salmon P. 1870. Majority opinion, \textit{Hepburn v.\
Griswold}, 75 U.S.\ (8 Wall.) 603.

\item Farrand, Max, ed. 1911. \textit{Records of the Federal Convention
of 1787}, 4 vols. New Haven: Yale University Press.

\item Field, Stephen J. 1871. Dissenting opinion, \textit{Knox v.\ Lee}
and \textit{Parker v.\ Davis}, 79 U.S.\ (12 Wall.) 457.

\item Field, Stephen J. 1884. Dissenting opinion, \textit{Juilliard
v.\ Greenman}, 110 U.S.\ 421.

\item Gray, Horace. 1884. Majority opinion, \textit{Juilliard
v.\ Greenman}, 110 U.S.\ 421. Text at
\texttt{law.cornell.edu/supremecourt/text/110/421}.

\item Hall, George J. and Thomas J. Sargent. 2014. ``Fiscal
Discriminations in Three Wars.'' \textit{Journal of Monetary Economics}
61: 148--166.

\item Kydland, Finn E. and Edward C. Prescott. 1977. ``Rules Rather
Than Discretion: The Inconsistency of Optimal Plans.''
\textit{Journal of Political Economy} 85(3): 473--491.

\item Madison, James. 1787. Notes of the Constitutional Convention,
August~16, 1787. Transcribed in Farrand (1911), vol.~2, pp.~304--310.
Online at \texttt{consource.org/document/james-madisons-notes-of-the%
-constitutional-convention-1787-8-16/}.

\item Marshall, John. 1819. \textit{McCulloch v.\ Maryland}, 17 U.S.\
(4 Wheat.) 316.

\item Marshall, John. 1830. \textit{Craig v.\ Missouri}, 29 U.S.\
(4 Pet.) 410.

\item Rogoff, Kenneth. 1985. ``The Optimal Degree of Commitment to
an Intermediate Monetary Target.'' \textit{Quarterly Journal of
Economics} 100(4): 1169--1189.

\item Strong, William. 1871. Majority opinion, \textit{Knox v.\ Lee}
and \textit{Parker v.\ Davis}, 79 U.S.\ (12 Wall.) 457. Text at
\texttt{tile.loc.gov/storage-services/service/ll/usrep/usrep079/%
usrep079457/usrep079457.pdf}.

\item United States. Statutes at Large, Vol.\ XII, Chap.\ XXXIII
(February 25, 1862). Legal Tender Act.

\item United States. Public Credit Act, March~18, 1869, 41st Cong.,
1st Sess., Chapter 1, 16 Stat.\ 1.

\end{description}

\end{document}
